\expandafter\def\csname HistoricalNoteText393\endcsname{Welche Mittel man auch dazu anwende, so wird man doch nicht darauf rechnen können, dass in völliger Strenge \textit{reine} Schwingungen um eine verticale Axe erzeugt werden. Schon aus diesem Grunde darf man von den Resultaten nicht die allerschärfste Übereinstimmung erwarten, worauf es jedoch hier auch gar nicht ankommt.}
\expandafter\def\csname HistoricalNoteText392\endcsname{Es ist nicht in allen Versuchen ganz gleich, da die Wirkung des Multiplicators sich mit einer, wie oben bemerkt ist, an sich nicht unveränderlichen Grösse verbindet. Auch hängt die Wirkung des Multiplicators selbst von dem mit der Temperatur veränderlichen Leitungsvermögen des Drahts mit ab.}
\SourceLeaf{392}Schwingungsgeschwindigkeit der Nadel am grössten ist, und der die entgegengesetzte Richtung annimmt, sobald die Nadel umkehrt: die Reaction dieses Stromes auf die Nadel besteht aber immer in einer Verminderung der Schwingungsgeschwindigkeit der letztern, und die Theorie ergibt, dass auch hievon eine Abnahme des Schwingungsbogens, sehr nahe in geometrischer Progression, die Folge sein muss. Da indessen die Intensität des inducirten Stromes auch durch den Widerstand, welchen die \textit{ganze} Kette darbietet, bedingt wird, so ist die Vergrösserung des logarithmischen Decrements, welche der Multiplicator hervorbringt, am stärksten, wenn die Kette gleich hinter diesem abgeschlossen ist; sie fällt desto kleiner aus, je grösser der hinzukommende Theil der Kette ist, und bei offener Kette findet gar kein Einfluss Statt, sondern das logarithmische Decrement ist dasselbe, als wenn der Multiplicator ganz weggenommen ist. Der Apparat, an welchem die obigen Beobachtungen angestellt sind, hat einen Multiplicator von 610 Umwindungen, und wenn derselbe, ohne weitern Zusatz, geschlossen ist, steigt das logarithmische Decrement auf etwa $0.02400$\HistoricalNote{392}, so dass der Schwingungsbogen schon in etwa 9 Minuten auf die Hälfte reducirt wird, während bei offener Kette und einem solchen Werthe des logarithmischen Decrements, wie die obigen Beobachtungen ergaben, etwa drei Stunden dazu erforderlich sind. Bei der vierpfündigen Nadel des magnetischen Observatorium bewirkt der Schluss eines aus 536 Windungen bestehenden Multiplicators ein logarithmisches Decrement von nahe gleicher Grösse; da aber jene Nadel eine Schwingungsdauer von $21''6$ hat, so kommt der Schwingungsbogen hier schon nach $4\tfrac12$ Minuten auf die Hälfte herab. Bei einer so bedeutenden Dämpfungskraft können, wenn die Nadel einmal beruhigt ist, falls nicht ausserordentliche Störungen von aussen oder ungewöhnlich starke schnelle Declinationsänderungen eintreten, gar keine erhebliche Schwingungen aufkommen, und ein solcher kräftiger Multiplicator gewährt daher, ausser seinen unzähligen andern Anwendungen, auch den wichtigen Nutzen, die Terminsbeobachtungen ungemein zu erleichtern, und alle andern Beruhigungsmittel entbehrlich zu machen.

In noch viel höherm Grade leistet aber diese Wirkung die oben [S. 372 d. B.] unter dem Namen eines Dämpfers erwähnte Vorrichtung. Der für das Magnetome\SourceLeaf{393}ter des magnetischen Observatorium angefertigte Dämpfer besteht in zwei länglich viereckigen kupfernen Rahmen, jeder 13 Pfund wiegend, welche in die hölzernen Rahmen der zwei Multiplicatorhälften wie eine Fütterung eingeschoben werden können. Da diese Vorrichtung in sehr vielen Fällen ungemein nützliche Dienste leistet, und auch schon mehrern vom Hrn. \textsc{Meyerstein} an auswärtige Beobachter gelieferten Magnetometern ein ähnlicher Hülfsapparat beigegeben ist, so werden einige denselben betreffende Bemerkungen hier nicht am unrechten Orte sein.

Zuvörderst wird jeder, welcher von einem solchen Apparat Gebrauch machen will, die Stärke seiner Dämpfungskraft in Zahlen kennen zu lernen wünschen. Man setzt zuerst die Nadel in sehr grosse Schwingungsbewegungen\HistoricalNote{393} und zeichnet die Elongationen, so bald sie innerhalb der Scale fallen, und so lange der Schwingungsbogen noch eine beträchtliche Grösse behält, auf. Sind $g$, $g'$, $g''$, $g'''$ u. s. w. die so hervorgehenden Amplituden in Scalentheilen, so erhält man in den Differenzen ihrer Logarithmen $\log g-\log g'$, $\log g'-\log g''$, $\log g''-\log g'''$ u. s. w. eben so viele Bestimmungen des logarithmischen Decrements. Zu einer Zeit, wo eine etwas beträchtliche Declinationsbewegung Statt findet, wird diese unter der Voraussetzung, dass sie während der Beobachtungen gleichförmig geschehe, eliminirt, wenn man sich der Formeln
\[
\log(g+g')-\log(g'+g''),\quad \log(g'+g'')-\log(g''+g''')\text{ u. s. w.}
\]
bedient, was, wie man leicht sieht, mit dem oben [S. 383 d. B.] angegebenen Verfahren auf Eins hinausläuft.

Einer der am 9. Januar 1838 angestellten Versuche gab z. B.
\begin{center}
\begin{tabular}{r|r|r|r}
Elongationen&Amplituden&Logarithmen&Differenzen\\
$134.0$&&&\\[-3pt]
&$1057.7$&$3.02436$&\\[-3pt]
$1191.7$&&&$0.31137$\\[-3pt]
&$516.4$&$2.71299$&\\[-3pt]
$675.3$&&&$0.30850$\\[-3pt]
&$253.8$&$2.40449$&\\[-3pt]
$929.1$&&&$0.31141$\\[-3pt]
&$123.9$&$2.09307$&\\[-3pt]
$805.2$&&&
\end{tabular}
\end{center}
\SourceLeaf{394}Also im Mittel das logarithmische Decrement $=0.31043$. Die andern Formeln geben
\[
\log\frac{1574.1}{770.2}=0.31043,\quad \log\frac{770.2}{377.7}=0.30945
\]
also im Mittel $0.30994$. Bei diesen Versuchen war der Multiplicator nicht geschlossen, oder es wirkte der Dämpfer allein. Aus andern Versuchen an demselben Tage, bei denen zugleich der Multiplicator geschlossen war, fand sich das logarithmische Decrement $=0.33570$.

Da Kupfer, wenn es nicht ganz rein ist, einen wenn auch nur sehr schwachen directen magnetischen Einfluss ausüben kann, so ist es nicht rathsam, den Dämpfer bei solchen Beobachtungen anzuwenden, die absolute Declinationsbestimmungen zum Zweck haben, ohne sich vorher überzeugt zu haben, dass ein merklicher Einfluss dieser Art nicht vorhanden ist. Man erfährt dies durch Beobachtungen des Standes der Nadel, abwechselnd mit und ohne Dämpfer, verbunden mit gleichzeitigen Beobachtungen des Standes einer in angemessener Entfernung befindlichen zweiten Nadel, um von den während der Beobachtungen Statt findenden Declinationsveränderungen Rechnung tragen zu können. In Ermangelung eines zweiten Magnetometers kann man diese Elimination, nur weniger zuverlässig, dadurch beschaffen, dass man die alternirenden Bestimmungen in nahe gleichen Zwischenzeiten macht, und jeden Stand ohne Dämpfer mit dem Mittel des vorhergehenden und folgenden Standes mit Dämpfer, und umgekehrt, vergleicht. Bei der Anfertigung des Dämpfers für das hiesige magnetische Observatorium hat Hr. \textsc{Meyerstein} die Vorsicht angewandt, sich nur ganz neuer Feilen zu bedienen; der fertige Dämpfer ist hernach eine beträchtliche Zeit in verdünnte Salzsäure, dann in Lauge gelegt, und zuletzt mit Wasser abgespült. Versuche der beschriebenen Art haben keinen merklichen Einfluss dieses Dämpfers auf den Stand der Nadel zu erkennen gegeben.

Was bei sich immer gleich bleibendem Schwingungsbogen strenge gültig sein würde, erleidet im Fall der Natur, wo der Schwingungsbogen fortwährend abnimmt, mehrere Modificationen, die hier noch etwas näher betrachtet zu werden verdienen, wäre es auch nur, um bestimmt beurtheilen zu können, unter welchen Umständen sie als unmerklich betrachtet werden dürfen.

Eine in geometrischer Progression erfolgende Abnahme des Schwingungsbogens setzt eine der Bewegung in jedem Augenblick in einfachem Verhältniss
\expandafter\def\csname HistoricalNoteText395\endcsname{\textit{Strenge} genommen gilt beides nur für unendlich kleine Schwingungen.}%
\SourceLeaf{395}ihrer Geschwindigkeit entgegenwirkende Kraft voraus\HistoricalNote{395}. Die allgemeine Gleichung für die Schwingungsbewegung hat daher, wenn wir die Grössen von der dritten Ordnung in Beziehung auf den Schwingungsbogen vernachlässigen, die Form
\[
0=\frac{\mathrm{d}\mathrm{d}x}{\mathrm{d}t^2}+nn(x-p)+2\epsilon.\frac{\mathrm{d}x}{\mathrm{d}t}
\]
wo $x$ den den Stand der Nadel für die Zeit $t$ bezeichnenden, $p$ den dem Ruhestande entsprechenden Scalentheil bedeuten, $nn$ und $2\epsilon$ hingegen die magnetische Directionskraft und jene retardirende Kraft, beide mit dem Trägheitsmoment der Nadel dividirt. Das vollständige Integral dieser Gleichung ist
\[
x=p+Ae^{-\epsilon t}\sin\BareRoot{(nn-\epsilon\epsilon)}.(t-B)
\]
wo $e$ die Basis der natürlichen Logarithmen, $A$ und $B$ die beiden durch die Integration eingeführten arbiträren Constanten bedeuten. Ohne die retardirende Kraft würde das Integral
\[
x=p+A\sin n(t-B)
\]
sein. Die Nadel macht also auch in jenem Fall wie in diesem periodische Oscillationen um den Punkt $p$, aber ein doppelter Unterschied findet dabei Statt. Theils ist im zweiten Fall die grösste Abweichung von der Mitte oder die halbe Amplitude constant $=A$, während sie im erstern in geometrischer Progression abnimmt, theils schreitet das Argument der periodischen Function im ersten Fall langsamer fort als im andern. Setzt man die Schwingungsdauer im zweiten Fall, wo sie allein von der magnetischen Directionskraft abhängt, $=T$, im ersten $=T'$, so hat man, $\pi$ in üblicher Bedeutung genommen,
\[
nT=\pi,\quad \BareRoot{(nn-\epsilon\epsilon)}.T'=\pi
\]
Wenn man also Kürze halber $n'$ anstatt $\BareRoot{(nn-\epsilon\epsilon)}$ schreibt, und einen Hülfswinkel $\varphi$ einführt, wonach
\[
\sin\varphi=\frac{\epsilon}{n},\quad \cos\varphi=\frac{n'}{n},\quad \operatorname{tang}\varphi=\frac{\epsilon}{n'}
\]
wird, wenn man ferner, wie oben, mit $\lambda$ das logarithmische Decrement und mit
\SourceLeaf{396}$m$ den Modulus des Systems bezeichnet, so erhält man
\[
\begin{gathered}
T'=\frac{T}{\cos\varphi}\\
\frac{\lambda}{m}=\epsilon T'=n'\operatorname{tang}\varphi.T'=\pi\operatorname{tang}\varphi
\end{gathered}
\]
folglich
\[
\operatorname{tang}\varphi=\frac{\lambda}{m\pi}=\frac{\lambda}{1.364376}
\]
Für $\lambda=0.02400$ und $T=42''18$ findet sich nach diesen Formeln $\varphi=1^0\,0'\,28''$ und $T'=42''18653$. Der blosse Schluss des Multiplicators bringt also nur eine geringe Vergrösserung der Schwingungsdauer hervor. Dagegen geben die oben beim Gebrauch des Dämpfers, allein, oder zugleich mit dem Multiplicator, gefundenen Zahlen, wenn man $T=20''60$ setzt,
\begin{center}
\begin{tabular}{r|r|r}
$\lambda=0.30994$&$\varphi=12^0\,47'\,54''$&$T'=21''12484$\\
$\lambda=0.33570$&$\varphi=13\quad49\quad22$&$T'=21.21439$
\end{tabular}
\end{center}

Die Beobachtungen stimmen mit dieser berechneten Vergrösserung der Schwingungsdauer so genau überein, als man nur von der geringen Anzahl von Schwingungen, auf die man sich dabei beschränken muss, erwarten kann.

In dem Fall abnehmender Schwingungsbögen sind die wahren Zeiten der Elongationen den aus correspondirenden Stellungen abgeleiteten nicht genau gleich, und bei so starken logarithmischen Decrementen, wie unter Anwendung eines Dämpfers Statt finden, wird dieser Unterschied ziemlich beträchtlich.

Da in dem oben gegebenen Integral offenbar $B$ die Zeit eines Durchganges durch den Ruhestand $p$ bedeutet, und es gleichgültig ist, von welchem Augenblick an die Zeit gezählt wird, so wollen wir grösserer Einfachheit wegen $B=0$ setzen. Unsere Formel wird so
\[
x=p+Ae^{-\epsilon t}\sin n't
\]
Der nächste Durchgang durch $p$, welcher auf den bei $t=0$ folgt, findet Statt bei $n't=180^0$, oder $t=T'$; die aus diesen correspondirenden Beobachtungen einfach abgeleitete Zeit der Elongation ist also $t=\tfrac12T'$, während der wirkliche Stillstand schon früher eintritt. Man hat nemlich für $\frac{\mathrm{d}x}{\mathrm{d}t}=0$,
\[
0=Ae^{-\epsilon t}(-\epsilon\sin n't+n'\cos n't)
\]
\SourceLeaf{397}Mithin
\[
\operatorname{cotang}n't=\frac{\epsilon}{n'}=\operatorname{tang}\varphi
\]
Daher der erste positive Werth von $n't=\tfrac12\pi-\varphi$, und $t=\tfrac12T'-\frac{\varphi T'}{\pi}$, oder in so fern $\varphi$ in Graden ausgedrückt ist,
\[
t=\tfrac12T'-\frac{\varphi}{180^0}.T'
\]

Offenbar findet eine gleiche Differenz bei der folgenden Elongation Statt. Aus den oben angegebenen Zahlen findet sie sich $=0''23$ für den fünfundzwanzigpfündigen Stab unter Anwendung des Multiplicators; $=1''50$ für das Magnetometer des M. O., wenn der Dämpfer allein, und $=1''63$, wenn Dämpfer und Multiplicator zugleich gebraucht werden.

Da die wirklichen Stillstände um eine constante, von der Grösse des Schwingungsbogens unabhängige, Zeit früher eintreten, als die aus aufeinanderfolgenden Durchgängen durch den Ruhestand $p$ geschlossenen Augenblicke, so kann man auch ohne Weiteres die letztern beibehalten, da es für den Gebrauch zur Bestimmung der Schwingungsdauer nur auf die \textit{Unterschiede} der Elongationszeiten ankommt. Nur muss man Sorge tragen, den Punkt $p$ selbst oder einen sehr nahe liegenden zum Beobachten zu wählen, und kleine Schwingungsbögen auch noch um so mehr ausschliessen, weil bei solchen in dem Fall starker logarithmischer Decremente schon eine geringe Abweichung vom richtigen $p$ einen merklichen Fehler erzeugen würde. Obgleich es nicht schwer ist, jener Bedingung Genüge zu leisten, so mag doch noch die allgemeine Formel für den Fehler der aus correspondirenden Durchgängen geschlossenen Elongationen hier Platz finden.

Es sei $u$ die halbe Zwischenzeit zwischen zwei correspondirenden Durchgängen eines Punkts $x$, und $\tfrac12T'-\delta$ das Mittel der Durchgangszeiten oder die daraus geschlossene Elongationszeit. Es sind also die Durchgangszeiten selbst $\tfrac12T'-u-\delta$ und $\tfrac12T'+u-\delta$, und wir haben folglich
\[
\begin{aligned}
x&=p+Ae^{-\epsilon(\tfrac12T'-u-\delta)}\sin n'(\tfrac12T'-u-\delta)\\
x&=p+Ae^{-\epsilon(\tfrac12T'+u-\delta)}\sin n'(\tfrac12T'+u-\delta)
\end{aligned}
\]
woraus, wegen $\tfrac12n'T'=\tfrac12\pi$, folgt
\[
e^{2\epsilon u}\cos n'(u+\delta)=\cos n'(u-\delta)
\]
\SourceLeaf{398}und mithin
\[
\operatorname{tang}n'\delta=\frac{e^{2\epsilon u}-1}{(e^{2\epsilon u}+1)\operatorname{tang}n'u}
\]

Für $u=\tfrac12T'$ gibt diese Formel $\delta=0$; dies ist der Fall, wo der Ruhestandspunkt $p$ selbst für die Durchgänge gewählt ist: hingegen entspricht die Annahme eines unendlich kleinen Werths dem wahren Stillstandspunkte, und die Formel gibt hier
\[
\operatorname{tang}n'\delta=\frac{\epsilon}{n'}=\operatorname{tang}\varphi
\]
Also $\delta=\frac{\varphi}{n'}=\frac{\varphi T'}{\pi}$, übereinstimmend mit dem oben gefundenen.

Endlich bedarf in dem Fall abnehmender Schwingungsbögen auch die Berechnung der auf den Ruhestand der Nadel bezüglichen Beobachtungen einer Modification, die freilich nur dann merklich wird, wenn die Schwingungen eine sehr starke Abnahme erleiden.

Die Stellungen der Nadel $x$, $x'$, welche zweien um eine Schwingungsdauer verschiedenen Zeiten $t$, $t+T'$ entsprechen, haben die Werthe
\[
\begin{aligned}
x&=p+Ae^{-\epsilon t}\sin n't\\
x'&=p+Ae^{-\epsilon t-\epsilon T'}\sin(n't+n'T')
\end{aligned}
\]
oder weil $n'T'=\pi$
\[
x'=p-Ae^{-\epsilon t-\epsilon T'}\sin n't
\]
oder wenn wir, wie oben, mit $\theta$ den Bruch bezeichnen, dessen briggischer Logarithme $-\lambda$, also der natürliche $\epsilon T'$ ist,
\[
x'=p-\theta Ae^{-\epsilon t}\sin n't
\]

Es erhellt also, dass man, um $p$ zu finden, nicht mehr das arithmetische Mittel zwischen $x$ und $x'$ nehmen darf, sondern die Differenz zwischen $x$ und $x$ in dem Verhältniss von 1 zu $\theta$ vertheilen muss, oder dass
\[
p=\frac{\theta x+x'}{1+\theta}=x+\frac{1}{1+\theta}.(x'-x)=x'-\frac{\theta}{1+\theta}.(x'-x)
\]
wird. Da übrigens bei diesen Beobachtungen die Differenz $x'-x$ immer sehr klein ist, so wird man zur Bequemlichkeit der Rechnung sich verstatten können, anstatt $\frac{\theta}{1+\theta}$ einen nahe kommenden durch kleine Zahlen auszudrückenden Bruch
\SourceLeaf{399}anzuwenden, z. B. kann man für $\lambda=0.30994$, anstatt des genaueren Werths $0.3288$ den genäherten $\tfrac13$ wählen.

Hiebei entsteht nun aber die Frage, für welchen Augenblick dieses Resultat als gültig zu betrachten ist. So wie in dem Falle, wo die Schwingungsbögen nur sehr langsam abnehmen, das einfache Mittel der Scalentheile als dem einfachen Mittel der Zeiten entsprechend angenommen wird, scheint nun zwar, dass bei ungleich vertheiltem Unterschied der Stände die Zwischenzeit in demselben Verhältniss zu theilen, also der Stand $x+\frac{x'-x}{1+\theta}$ als für $t+\frac{T'}{1+\theta}$ gültig anzusehen sei: allein dies ist theoretisch nicht richtig, und es scheint eine genauere Erörterung, wenn auch in gewöhnlichen Fällen praktisch ganz unerheblich, doch in theoretischer Beziehung hier noch eines Platzes nicht unwürdig zu sein.

Offenbar kommt der Gültigkeitsaugenblick nur in so fern in Frage, als man sich nicht erlauben will, die magnetische Declination in der Zwischenzeit zwischen den beiden Aufzeichnungen als constant zu betrachten: aber als sich gleichförmig während dieser Zwischenzeit ändernd wird man sie immer betrachten können, und müssen, wenn der Rechnung eine bestimmte Unterlage gegeben werden soll\quad In diesem Falle hat also unsere Fundamentalgleichung die Form
\[
0=\frac{\mathrm{d}\mathrm{d}x}{\mathrm{d}t^2}+nn(x-p-\alpha t)+2\epsilon.\frac{\mathrm{d}x}{\mathrm{d}t}
\]
deren vollständiges Integral ist
\[
x=p-\frac{2\alpha\epsilon}{nn}+\alpha t+Ae^{-\epsilon t}\sin n'(t-B)
\]
wenn, wie oben, $n'$ für $\BareRoot{(nn-\epsilon\epsilon)}$ gesetzt wird. Wenn also $x'$ den Stand für die Zeit $t+T'$ ausdrückt, so wird, $\theta$ in voriger Bedeutung genommen,
\[
x'=p-\frac{2\alpha\epsilon}{nn}+\alpha t+\alpha T'-\theta Ae^{-\epsilon t}\sin n'(t-B)
\]
und folglich
\[
\frac{\theta x+x'}{1+\theta}=p-\frac{2\alpha\epsilon}{nn}+\alpha t+\frac{\alpha T'}{1+\theta}
\]
welches Resultat demnach der Ruhestand für die Zeit
\[
t+\frac{T'}{1+\theta}-\frac{2\epsilon}{nn}=t+\frac{T'}{1+\theta}-\frac{\sin2\varphi.T'}{\pi}
\]
ist. Für $\lambda=0.30994$ ist also die Zeit, wofür das nach obiger Vorschrift berechnete Resultat gilt $=t+0.5337T'$, für $\lambda=0.33570$ hingegen $=t+0.5365T'$.
\SourceLeaf{400}Man sieht also, dass selbst bei einer so starken Dämpfung der Augenblick der Gültigkeit von dem einfachen Mittel der Zeiten nur wenig verschieden ist.

Bei allem, was bisher entwickelt ist, liegt die Voraussetzung zum Grunde, dass $\epsilon$ kleiner sei als $n$; im entgegengesetzten Fall nimmt das Integral der Fundamentalgleichung eine andere Form an. Man erhält nemlich anstatt des Gliedes $Ae^{-\epsilon t}\sin\BareRoot{(nn-\epsilon\epsilon)}.(t-B)$, in dem Fall, wo $\epsilon$ grösser ist als $n$, zwei Glieder von der Form
\[
Ae^{-(\epsilon+\BareRoot{(\epsilon\epsilon-nn)})t}+Be^{-(\epsilon-\BareRoot{(\epsilon\epsilon-nn)})t}
\]
und in dem Fall, wo $\epsilon=n$ ist, von dieser
\[
(A+Bt)e^{-\epsilon t}
\]

In beiden Fällen findet also in der Bewegung gar nichts periodisches mehr Statt, sondern der Stand nähert sich asymptotisch dem Ruhestande. Für unsern Dämpfer ist $\frac{\epsilon}{n}=0.22152$, und es müsste also ein mehr als $4\tfrac12$ mal stärker wirkender Dämpfer angewandt werden, um solchen Erfolg hervorzubringen. Offenbar aber würde es dazu nicht hinreichend sein, die Metallmenge nur in demselben Verhältniss zu vergrössern, in sofern diese Vergrösserung nach aussen angebracht werden müsste, und die äussern Schichten des Metallrahmens vergleichungsweise weniger zur Inductionswirkung beitragen als die innern. Allein es würde nicht einmal anzurathen sein, eine Dämpfung von einer solchen Stärke anzuwenden, dass die Bewegung aufhörte periodisch zu sein, theils weil, sobald $\epsilon$ den Grenzwerth $n$ überschreitet, die Annäherung zu dem Ruhestand wieder langsamer geschieht, theils weil man dann den wesentlichen Vortheil verlöre, aus zwei beliebigen, um $T'$ von einander entfernten, Aufzeichnungen den Ruhestand auf eine bequeme Art berechnen zu können.

\clearpage
\fancyhead[L]{\small BEOBACHTUNG VON ABLENKUNGEN.}
\SourceLeaf{401}
\pdfbookmark[0]{Über ein Mittel die Beobachtung von Ablenkungen zu erleichtern. 1839}{ablenkungen-1839}
\begin{center}
{\large ÜBER EIN MITTEL\par\medskip DIE BEOBACHTUNG VON ABLENKUNGEN ZU ERLEICHTERN.}
\par\medskip
Resultate aus den Beobachtungen des magnetischen Vereins. 1839. II.
\end{center}
\par\medskip\Needspace{4\baselineskip}\pdfbookmark[1]{1.}{ablenkungen-art-1}\centerline{\textbf{1.}}\nobreak\smallskip
Wenn zu der erdmagnetischen Kraft noch eine andere auf die Nadel eines Magnetometers stetig, aber in einer gegen den magnetischen Meridian geneigten Richtung wirkende Kraft hinzutritt, so erhält die Nadel eine veränderte Gleichgewichtsstellung, und die Grösse der Ablenkung kann zur Abmessung der Zusatzkraft dienen. Zur Messbarkeit der Ablenkung ist aber erforderlich, dass nicht nur die neue Gleichgewichtsstellung noch innerhalb der Scale liege, sondern auch, insofern man nicht den völligen Ruhezustand der Nadel abwarten kann oder will, dass die noch Statt findenden Schwingungen die Grenzen der Scale nicht überschreiten. War die Nadel, so lange der Erdmagnetismus allein auf sie wirkte, in Ruhe, und setzt man die Zusatzkraft auf einmal in volle Wirkung, so fängt jene eine Schwingung an, deren Mitte die neue Gleichgewichtsstellung ist, während die vorige Stellung den einen Elongationspunkt bildet, und der zweite eben soweit von der Mitte auf der entgegengesetzten Seite hinausfällt. Liegt nun die neue Gleichgewichtsstellung zwar innerhalb, aber doch nahe an der Grenze der Scale, so würde man bei der langsamen Abnahme des Schwingungsbogens ohne Anwendung künstlicher Hülfsmittel auf diese Art erst lange zu warten haben, bis die Bestimmung jenes Punkts möglich würde. Dadurch würde aber in allen Fällen schon wegen der stündlichen Veränderung der Declination, die Zuverlässigkeit und Brauchbarkeit der Bestimmung sehr vermindert, und fast ganz vereitelt
